% COURSE-ASSISTANT SOURCE — FALL 2025 ARCHIVE
% This material is not authoritative for Fall 2026; current-course material
% takes priority. See LN04_figure_notes.md for figures not stored here and
% LN04_errata.md for a correction to the human-capital growth-accounting slide.
\documentclass[10pt,handout]{beamer}
\input{EC2211_preamble.txt}
\title{Lecture Notes 4\newline TFP and Long-Term Economic Growth}


\begin{document}
\maketitle


\begin{frame}{Contents and literature}
    \begin{itemize}
        \item Towards endogenous growth
        \item The Romer model
        \item Institutions and growth
        \item Growth at the frontier
    \end{itemize}
    \bigskip

    Literature: 
    \small{
    \begin{tightitemize}
        \item Jones (2024), chapter 6
        \item The Royal Swedish Academy of Sciences (2018), \href{https://www.nobelprize.org/uploads/2018/10/popular-economicsciencesprize2018.pdf}{"Integrating nature and knowledge into economics"}, in particular pages 1-4
        \item Royal Swedish Academy of Sciences (2024) \href{https://www.nobelprize.org/uploads/2024/10/popular-economicsciencesprize2024.pdf}{"They provided an explanation for why some countries are rich and others poor"}
        \item Gordon (2014), \href{https://www.nber.org/system/files/working_papers/w19895/w19895.pdf}{"The demise of U.S. economic growth"}, sections 1-3.3
        \item Jones (2023), \href{https://www.kansascityfed.org/Jackson\%20Hole/documents/9771/GrowthOutlookPanel.pdf}{"The outlook for long-term economic growth"}
        \item Draghi (2024) \href{https://commission.europa.eu/document/download/97e481fd-2dc3-412d-be4c-f152a8232961_en?filename=The\%20future\%20of\%20European\%20competitiveness\%20_\%20A\%20competitiveness\%20strategy\%20for\%20Europe.pdf}{"The future of European competitiveness"}, pages 5-19
    \end{tightitemize}
    }
\end{frame}



%%%%%%%%%%%%%%%%%%%%%%%%%%%%%%%%%%%%%%%%%%%%%%%%%%%%%%%%%%%%%%%%%%%%%%%%%%%%%%%%%%%%%%%%%%%%%%
%%%%%%%%%%%%%%%%%%%%%%%%%%%%%%%%%%%%%%%%%%%%%%%%%%%%%%%%%%%%%%%%%%%%%%%%%%%%%%%%%%%%%%%%%%%%%%
\section{Towards endogenous growth}


\begin{frame}{Towards endogenous growth}
\begin{itemize}
    \item The Solow model offers valuable insights into cross-country differences in income per capita but fails to explain sustained growth\medskip
    \item Next: endogenous growth
\end{itemize}
\end{frame}


\begin{frame}{The AK "model"}
    What happens when we relax the assumption of diminishing marginal product of capital? \smallskip

    \begin{equation}\label{eq35}
        Y_t=AK_t
    \end{equation}\smallskip

    Capital accumulation:
    \begin{equation}\label{eq36}
        K_{t+1} = (1-\delta)K_t + I_t = (1-\delta) K_t + s AK_t
    \end{equation}\smallskip

    Equations \eqref{eq35} and \eqref{eq36} imply that:
    \begin{equation}\label{eq37}
         \frac{Y_{t+1}}{Y_{t}}=\frac{K_{t+1}}{K_{t}}= 1 - \delta + sA
    \end{equation}
\end{frame}


\begin{frame}{The AK "model"}
    Since
    \begin{equation*}
        \gamma \equiv \frac{Y_{t+1}-Y_t}{Y_{t}}= \frac{Y_{t+1}}{Y_{t}}-1,
    \end{equation*}
    the growth rate of GDP is:
    \begin{equation}\label{eq38}
        \gamma = sA - \delta
    \end{equation}

    An assumption of constant marginal product of capital thus generates endogenous growth! \medskip

    Recall that a higher savings rate $s$ did \emph{not} generate higher long-run growth in the Solow model\medskip

    We see here that higher savings (more investments) would generate higher long-run growth if the marginal product of capital were not diminishing
\end{frame}



%%%%%%%%%%%%%%%%%%%%%%%%%%%%%%%%%%%%%%%%%%%%%%%%%%%%%%%%%%%%%%%%%%%%%%%%%%%%%%%%%%%%%%%%%%%%%%
%%%%%%%%%%%%%%%%%%%%%%%%%%%%%%%%%%%%%%%%%%%%%%%%%%%%%%%%%%%%%%%%%%%%%%%%%%%%%%%%%%%%%%%%%%%%%%
\section{The Romer model}


\begin{frame}{The Romer model}
\begin{columns}
    \column{0.6\textwidth}
    \begin{itemize}
        \item Romer (1990), "Endogenous technological change"\medskip

        \item \href{https://www.nobelprize.org/prizes/economic-sciences/2018/romer/facts/}{Nobel prize 2018}: \small{"for integrating technological innovations into long-run macroeconomic analysis"}
    \end{itemize}

    \column{0.4\textwidth}
    \begin{figure}
        \centering
        \includegraphics[width=0.95\linewidth]{romer.jpg}
    \end{figure}
    \centering \small{Paul Romer 1955-}
\end{columns}
\end{frame}


\begin{frame}{The Romer model}
\begin{itemize}
    \item Distinguish between objects (rivalrous) and ideas (non-rivalrous)\medskip
    \item Workers can either produce consumption goods or new ideas\medskip
    \item A single idea can be used by many people simultaneously\medskip
    \item Non-rivalry generates increasing returns\medskip
    \item Omit physical capital for simplicity
\end{itemize}
\end{frame}


\begin{frame}{}
    \begin{figure}
        \centering
        \includegraphics[width=1\linewidth]{Nonrivalry goods.jpg}
    \end{figure}
    \flushright \tiny{From Royal Academy of Sciences (2018)}
\end{frame}


\begin{frame}{Notation}
\begin{itemize}
    \item[] $Y_t$: output of the consumption good
    \item[] $A_t$: the existing stock of knowledge
    \item[] $L_t$: the population
    \item[] $n$: population growth rate
    \item[] $L_{yt}$: number of workers employed in goods production
    \item[] $L_{at}$: number of researchers (producing ideas)
    \item[] $\gamma$: parameter in the production function 
    \item[] $\lambda$: fraction of labor force employed in research
    \item[] $\zeta$: productivity parameter
    \item[] $A_0$: the initial stock of knowledge
\end{itemize}
\end{frame}


\begin{frame}{The Romer model}
    Production of the consumption good:
    \begin{equation}\label{eq39}
        Y_t = A_t^\gamma L_{yt}
    \end{equation}

    Production of new ideas:
    \begin{equation}\label{eq40}
        \Delta A_{t+1} = \zeta L_{at}
    \end{equation}

    Resource constraint for labor:
    \begin{equation}\label{eq41}
        L_{yt} + L_{at} = L_t
    \end{equation}
    and 
        \[L_{t+1}=(1+n)L_t
    \] 
    
    Labor is allocated according to:
    \begin{equation}\label{eq42}
        L_{at} = \lambda L_t
    \end{equation}
    \begin{equation}\label{eq43}
        L_{yt} = (1-\lambda)L_t
    \end{equation}
\end{frame}


\begin{frame}{Solving the model}
    Dividing \eqref{eq39} by $L$ and using (\ref{eq43}) gives output per capita: 
    \begin{equation}\label{eq44}
        y_t \equiv \frac{Y_t}{L_t} = A_t^\gamma (1-\lambda)
    \end{equation}

    Dividing \eqref{eq40} by $A_t$ gives the growth rate of TFP:
    \begin{equation*}
        g_{At} \equiv \frac{\Delta A_{t+1}}{A_t} = \zeta \frac{L_{at}}{A_t}
    \end{equation*}
    or (assuming that the growth rate is positive)
    \begin{equation}\label{eq45}
        A_t = \frac{\zeta}{g_{At}}L_{at}
    \end{equation}
\end{frame}


\begin{frame}{The balanced growth path}
\begin{itemize}
    \item Is there a long-run equilibrium ("balanced growth path") where all variables grow at constant rates? \smallskip

    \item If so, $g_{At}$ must be constant (since we assume that $L$, $L_y$ and $L_a$ grow at constant rates)
    \begin{tightitemize}
        \item[] Equation \eqref{eq45} then shows that the stock of knowledge is proportional to the number of researchers, $L_{at}$
    \end{tightitemize} \smallskip

    \item Since $L_{at}$ in turn is a constant fraction of the population, this means that $A$ and $L$ must grow at the same rate on the balanced growth path
    \begin{tightitemize}
        \item[] That is, $g_A=n$ in the long run
    \end{tightitemize}
\end{itemize}
\end{frame}


\begin{frame}{Interpreting the Romer model}
\begin{itemize}
    \item Equation \eqref{eq44} shows that output per capita will grow as long as the stock of knowledge, $A_t$ increases\smallskip

    \item The model moreover shows that the stock of knowledge increases at the same rate as the population\smallskip

    \item And from the production function we see that the growth rate of output per capita is $\gamma n$\smallskip

    \item Note the key difference to the Solow model where an investment generates more capital, but where capital is rivalry -- it is "less useful" when shared by many workers\smallskip

    \item In the Romer model, investment generates more knowledge which is not rivalry; all workers become more productive irrespective of how many they are
\end{itemize}
\end{frame}


\begin{frame}{A permanent rise in research productivity (higher $\zeta$)}
    \pause
    \begin{figure}
        \includegraphics[scale=0.35]{MACRO6_FIG06.03.jpg}
    \end{figure}
    \flushright \tiny Figure 6.3 in Jones (2024). Recall that his $\overline{z}$ is our $\zeta$.
\end{frame}


\begin{frame}{A permanent rise in the research share (higher $\lambda$)}
    \pause
    \begin{figure}
        \includegraphics[scale=0.35]{MACRO6_FIG06.04.jpg}
    \end{figure}
    \flushright \tiny Figure 6.4 in Jones (2024). Recall that his $\overline{l}$ is our $\lambda$.
\end{frame}



%%%%%%%%%%%%%%%%%%%%%%%%%%%%%%%%%%%%%%%%%%%%%%%%%%%%%%%%%%%%%%%%%%%%%%%%%%%%%%%%%%%%%%%%%%%%%%
%%%%%%%%%%%%%%%%%%%%%%%%%%%%%%%%%%%%%%%%%%%%%%%%%%%%%%%%%%%%%%%%%%%%%%%%%%%%%%%%%%%%%%%%%%%%%%
\section{Summing up:\newline Malthus $\rightarrow$ Solow $\rightarrow$ Romer}


\begin{frame}{Malthus $\rightarrow$ Solow $\rightarrow$ Romer}
\begin{itemize}
    \item What is important in production? \newline
            Land (Malthus) $\rightarrow$ Physical capital (Solow) $\rightarrow$ Ideas (Romer) \bigskip
    \item Population growth rates
\end{itemize}  
\end{frame}


\begin{frame}{Total fertility rates have trended down}
\begin{columns}
    \column{0.5\linewidth}
    \begin{figure}
        \centering
        \includegraphics[width=1\linewidth]{Figs/TFR_a.png}
    \end{figure}
    \column{0.5\linewidth}
    \begin{figure}
        \centering
        \includegraphics[width=1\linewidth]{Figs/TFR_b.png}
    \end{figure}
\end{columns}
\flushright \tiny Total fertility rates (the number of children that would be born to a woman if she were to live to the end of her childbearing years and bear children in accordance with age-specific fertility rates of the specified year). Source: \href{https://data360.worldbank.org/en/indicator/WB_WDI_SP_DYN_TFRT_IN?view=datatable&recentYear=false}{World Bank}
\end{frame}


\begin{frame}{Malthus $\rightarrow$ Solow $\rightarrow$ Romer}
\begin{itemize}
    \item What is important in production? \newline
            Land (Malthus) $\rightarrow$ Physical capital (Solow) $\rightarrow$ Ideas (Romer) \bigskip
    \item Population growth rates
    \begin{tightitemize}
        \item Malthus: Population growth is \emph{endogenous}. (Do falling birth rates indicate that productivity has fallen?)
        \item Solow: A lower population growth rate means that we do not have to build as much capital for our children
        \item Romer: A lower population growth rate means that there will be fewer new ideas
    \end{tightitemize}
\end{itemize}  
\end{frame}



%%%%%%%%%%%%%%%%%%%%%%%%%%%%%%%%%%%%%%%%%%%%%%%%%%%%%%%%%%%%%%%%%%%%%%%%%%%%%%%%%%%%%%%%%%%%%%
%%%%%%%%%%%%%%%%%%%%%%%%%%%%%%%%%%%%%%%%%%%%%%%%%%%%%%%%%%%%%%%%%%%%%%%%%%%%%%%%%%%%%%%%%%%%%%
\section{Institutions and growth}


\begin{frame}{The role of institutions}
    \begin{quote}
        "The factors we have listed (innovation, economies of scale, education, capital accumulation, etc.) are not causes of growth; they are growth."
    \end{quote}\bigskip
    \begin{flushright}-North, D.C. and Thomas, R.P. (1973),\linebreak \textit{The Rise of the Western World: A New Economic History}, Cambridge University Press, Cambridge, UK.\end{flushright}
\end{frame}


\begin{frame}{Proximate versus fundamental causes of growth}
\begin{itemize}
    \item If growth is about factor accumulation and innovation, why do not poor countries invest and innovate?\medskip

    \item Why do countries provide different incentives to firms and households?\medskip

    \item Proximate causes of growth: factor accumulation and investment in R\&D\medskip

    \item Fundamental causes of growth: luck? geography? culture? institutions?\medskip
\end{itemize}
\end{frame}


\begin{frame}{Institutions and growth (Acemoglu, Johnson and Robinson)}
    \begin{itemize}
        \item Institutions
        \begin{tightitemize}
            \item Humanly devised set of rules
            \item Place constraints on individual behaviour
        \end{tightitemize}\medskip

    \item Economic institutions
    \begin{tightitemize}
        \item Property rights
        \item The presence and functioning of markets
        \item Contractual opportunities available to firms and households
        \item Entry barriers
    \end{tightitemize}\medskip

    \item Political institutions
    \begin{tightitemize} 
        \item Polity and electoral laws
        \item Checks and balances on political leaders
    \end{tightitemize}
\end{itemize}
\end{frame}


\begin{frame}{Democracy and income}
\begin{itemize}
    \item Is democracy conducive to high income?\medskip

    \item All the old democracies of the world are rich\medskip 

    \item Growth performance of autocracies highly heterogeneous
    \begin{tightitemize}
        \item Dictatorships excel or fall behind
        \item Many of the East Asian growth miracles started out autocratic
    \end{tightitemize}\medskip

    \item Complex relationship between polity and income!
\end{itemize}
\end{frame}



%%%%%%%%%%%%%%%%%%%%%%%%%%%%%%%%%%%%%%%%%%%%%%%%%%%%%%%%%%%%%%%%%%%%%%%%%%%%%%%%%%%%%%%%%%%%%%
%%%%%%%%%%%%%%%%%%%%%%%%%%%%%%%%%%%%%%%%%%%%%%%%%%%%%%%%%%%%%%%%%%%%%%%%%%%%%%%%%%%%%%%%%%%%%%
\section{Growth at the frontier}


\begin{frame}{Economic growth: summing up}
    Sources of temporary and/or long-run growth:\smallskip
    \begin{itemize}
        \item Physical capital\smallskip
        \item Labor, hours worked\smallskip
        \item Human capital\smallskip
        \item Productivity, technology
        \begin{tightitemize}
            \item Catching up with the frontier
            \item Efficient use of resources and people
            \item New ideas, research
        \end{tightitemize}
    \end{itemize}
\end{frame}


\begin{frame}{Growth accounting: Can we quantify the contribution of different sources?}
    Consider our Cobb-Douglas production function:
    \begin{equation*}
        Y_t = A_t K_t^\alpha L_t^{1-\alpha}
    \end{equation*} \medskip

    Apply the rules for growth rates:
    \begin{equation*}
        g_{Y_t} = g_{A_t} + \alpha g_{K_t} + (1-\alpha) g_{L_t}
    \end{equation*} \medskip

    In per capita terms, production is $y_t = A_t k_t^\alpha$, so
    \begin{equation*}
        g_{y_t} = g_{A_t} + \alpha g_{k_t}
    \end{equation*} \medskip
\end{frame}


\begin{frame}{Growth accounting for the USA}
\begin{figure}
    \centering
    \includegraphics[width=1\linewidth]{MACRO6_Table06.02.jpg}
\end{figure}
\flushright \tiny{Table 6.2 in Jones}
\end{frame}


\begin{frame}{Growth accounting with human capital}
    Suppose that the production function is:
    \begin{equation*}
        Y_t = A_t K_t^\alpha (h_tL_t)^{1-\alpha}
    \end{equation*} 
    where $h$ is human capital per worker\bigskip

    In per capita terms:
    \begin{equation*}
        y_t = A_t k_t^\alpha h_t^{1-\alpha} 
    \end{equation*} \medskip

    Then:
    \begin{equation*}
        g_{y_t} = g_{A_t} + \alpha g_{K_t} + (1-\alpha) g_{h_t}
    \end{equation*} 
\end{frame}


\begin{frame}{Growth accounting with human capital}
    Can we measure human capital or do we have to put it in the unexplained TFP term? \medskip

    Human capital is often measured by the average number of years of schooling in the adult population. This is then converted to a human capital index based on empirical data on how schooling affects wages.\medskip

    The \href{https://www.rug.nl/ggdc/productivity/pwt/?lang=en}{Penn World Tables} report such an index:\smallskip

    \begin{table}
    \small{
        \centering
        \begin{tabular}{lccc}
             & Years of schooling &  & "Human capital" \\ \hline
            USA  & 13.6 & & 3.75 \\
            Sweden & 12.3 & & 3.44 \\
            Mexico & 9.2 & & 2.78 \\
            Malawi & 5.7 & & 0.71
        \end{tabular}
        }
    \end{table}
\end{frame}


\begin{frame}{The outlook for long-term economic growth}
    Growth in the U.S. has been surprisingly stable at 2\% per year for the past 150 years\bigskip

    Gordon (2014) and Jones (2023) note that much of this growth is a result of
    \begin{itemize}
        \item Rising educational attainment
        \item Rising investment in research
    \end{itemize}\medskip

    I would add globalization (with offshoring, restructuring etc) to the list 
\end{frame}


\begin{frame}{The outlook for long-term economic growth}
\begin{figure}
    \centering
    \includegraphics[width=0.9\linewidth]{JonesGrowthAccounting.png}
\end{figure}
\flushright \tiny Figure 5 in Jones (2023)
\end{frame}


\begin{frame}{The outlook for long-term economic growth}
    Potential tailwinds:
    \begin{itemize}
        \item Better utilization of resources (use of talents)
        \item More frontier contribution from China and India
        \item AI
    \end{itemize}\medskip

    Headwinds:
    \begin{itemize}
        \item Educational attainment cannot/will not continue to increase
        \item Research intensity cannot/will not continue to increase
        \item Population growth will slow or turn negative
        \item Geopolitics, end of globalization
    \end{itemize}
\end{frame}


\begin{frame}{What we did (growth part of the course)}
\begin{itemize}
    \item Growth facts\medskip
    \item Growth accounting\medskip
    \item From Malthus to Solow to Romer \medskip
    \item Institutions
\end{itemize}
\end{frame}


\end{document}
